\section {Assunzione impegni}
Con la sottoscrizione del presente documento, ciascun componente del gruppo \textbf{Git Happens} si impegna formalmente a dedicare allo svolgimento del progetto un monte orario complessivo di \textbf{93 ore produttive ciascuno}, per un totale di \textbf{651 ore produttive} riferite all'intero team di 7 persone.
\\
Ciascun membro dichiara di aver compreso il carico di lavoro complessivo e si impegna a svolgere con continuità e responsabilità le attività previste, rispettando i ruoli, le scadenze concordate ed il \textit{Way of Working} definito dal gruppo.

\section{Analisi dei ruoli}
Tutti i ruoli tipici del ciclo di vita del software saranno coperti a rotazione da ogni membro del team per garantire una formazione omogenea ed un contributo equamente bilanciato:
\begin{itemize}
    \item \textbf{Responsabile:} Figura di riferimento per l'intero team; si occupa della pianificazione temporale delle attività, della gestione delle risorse, del coordinamento del lavoro e del rispetto dei vincoli stabiliti.
    \item \textbf{Amministratore:} Gestisce i processi, la configurazione degli ambienti di lavoro (GitHub, progetti Kanban, pipeline CI/CD) e garantisce l'efficienza degli strumenti operativi.
    \item \textbf{Analista:} Si occupa dell'analisi approfondita dei requisiti, raccogliendo le esigenze degli stakeholder e traducendole in specifiche dettagliate e casi d'uso.
    \item \textbf{Progettista:} Definisce l'architettura complessiva del sistema, identificando moduli, componenti e interazioni a partire dalle specifiche definite dagli analisti.
    \item \textbf{Programmatore:} Responsabile dell'implementazione e della scrittura del codice sorgente, dei test di unità e del refactoring in conformità alle specifiche.
    \item \textbf{Verificatore:} Cura il controllo qualità formale su codice e documentazione, verificando la conformità ai requisiti e l'assenza di difetti.
\end{itemize}